Se deben describir los recursos que se usan en la investigación:

\begin{enumerate}
	\item Equipo de Computo:
	\begin{itemize}
		\item Nombre: LAPTOP-EHFPNJP6.
		\item Procesador	11th Gen Intel(R) Core(TM) i5-11400H @ 2.70GHz   2.69 GHz
		\item RAM instalada	16.0 GB (15.7 GB usable)
		\item Tipo de sistema	Sistema operativo de 64 bits, procesador basado en x64
	\end{itemize}

    \item Google Colaboratory (Colab)
    \begin{itemize}
        \item Plataforma de computación en la nube proporcionada por Google que permite ejecutar código Python en notebooks interactivos. Se emplea para el entrenamiento de modelos de deep learning, procesamiento de señales EEG y visualización de resultados. Facilita el acceso a recursos como GPU y TPU sin necesidad de instalación local.
    \end{itemize}
    
    \item BNCI Horizon 2020
    \begin{itemize}
    \item Repositorio público de conjuntos de datos EEG utilizado ampliamente en investigación en interfaces cerebro-computadora (BCI). Contiene datos anotados para tareas como imaginación motora, estados cognitivos y estímulos visuales o auditivos. 
    \item Enlace: \url{https://bnci-horizon-2020.eu/database/data-sets}
    \end{itemize}

    \item PhysioNet
    \begin{itemize}
    \item Plataforma que ofrece acceso abierto a datos fisiológicos, incluyendo EEG. El conjunto de datos \textit{Auditory EEG Dataset} contiene registros de respuesta cerebral a estímulos auditivos, útil para el análisis de procesamiento sensorial y cognitivo.
    \item Enlace: \url{https://physionet.org/content/auditory-eeg/1.0.0/}
\end{itemize}

\end{enumerate}
